% from ega0-1-fr.tex
\oldpage[0\textsubscript{I}]{31}

Si la catégorie $\mathbf{K}$ admet des limites projectives (généralisées), et
si $\mathfrak{B}$, $\mathfrak{B}'$ sont des bases des topologies de $X$ et
$Y$ respectivement, pour définir un $\psi$-morphisme $u$ de \emph{faisceaux},
on peut se borner à se donner les $u_{U,V}$ pour
$U\in\mathfrak{B}$, $V\in\mathfrak{B}'$ et $\psi(U)\subset V$, vérifiant les
conditions de compatibilité (3.5.1.1) pour $U$, $U'$ dans $\mathfrak{B}$ et
$V$, $V'$ dans $\mathfrak{B}'$~; il suffit en effet de définir $u_W$, pour
tout ouvert $W\subset Y$, comme limite projective des $u_{U,V}$ pour
$V\in\mathfrak{B}'$ et $V\subset W$, $U\in\mathfrak{B}$ et
$\psi(U)\subset V$.

Lorsque la catégorie $\mathbf{K}$ admet des limites inductives, on a, pour
tout $x\in X$, un morphisme
\[
  \mathcal{G}(V)\to\mathcal{F}(\psi^{-1}(V))\to\mathcal{F}_x
\]
pour tout voisinage ouvert $V$ de $\psi(x)$ dans $Y$, et ces morphismes
forment un système inductif qui donne par passage à la limite un morphisme
\[
  \mathcal{G}_{\psi(x)}\to\mathcal{F}_x.
\]
\end{env}

\begin{env}[3.5.2]
\label{0.3.5.2-fr}
Sous les hypothèses de (3.4.3), soient $\mathcal{F}$, $\mathcal{G}$,
$\mathcal{H}$ des préfaisceaux à valeurs dans $\mathbf{K}$ sur $X$, $Y$, $Z$
respectivement, et soient
$u:\mathcal{G}\to\psi_*(\mathcal{F})$,
$v:\mathcal{H}\to\psi'_*(\mathcal{G})$ un $\psi$-morphisme et un
$\psi'$-morphisme respectivement. On en déduit un $\psi''$-morphisme
\[
  w:\mathcal{H}\xrightarrow{v}\psi'_*(\mathcal{G})
  \xrightarrow{\psi'_*(u)}
  \psi'_*(\psi_*(\mathcal{F}))=\psi''_*(\mathcal{F}),
\]
que l'on appelle, par définition, le \emph{composé} de $u$ et de $v$. On peut
donc considérer les couples $(X,\mathcal{F})$ formés d'un espace topologique
$X$ et d'un préfaisceau $\mathcal{F}$ sur $X$ (à valeurs dans $\mathbf{K}$)
comme formant une \emph{catégorie}, les morphismes étant les couples
$(\psi,\theta):(X,\mathcal{F})\to(Y,\mathcal{G})$ formés d'une application
continue $\psi:X\to Y$ et d'un $\psi$-morphisme
$\theta:\mathcal{G}\to\mathcal{F}$.
\end{env}

\begin{env}[3.5.3]
\label{0.3.5.3-fr}
Soient $\psi:X\to Y$ une application continue, $\mathcal{G}$ un
\emph{préfaisceau} sur $Y$ à valeurs dans $\mathbf{K}$. Nous appellerons
\emph{image réciproque de $\mathcal{G}$ par $\psi$} un couple
$(\mathcal{G}',\rho)$, où $\mathcal{G}'$ est un \emph{faisceau} sur $X$ à
valeurs dans $\mathbf{K}$, et
$\rho:\mathcal{G}\to\mathcal{G}'$ un $\psi$-morphisme (autrement dit un
homomorphisme $\mathcal{G}\to\psi_*(\mathcal{G}'))$ tels que, pour tout
\emph{faisceau} $\mathcal{F}$ sur $X$ à valeurs dans $\mathbf{K}$,
l'application
\begin{equation*}
  \operatorname{Hom}_X(\mathcal{G}',\mathcal{F})
  \longrightarrow \operatorname{Hom}_\psi(\mathcal{G},\mathcal{F})
  =\operatorname{Hom}_Y(\mathcal{G},\psi_*(\mathcal{F}))
\tag{3.5.3.1}\label{0.3.5.3.1-fr}
\end{equation*}
transformant $v$ en $\psi_*(v)\circ\rho$, soit une \emph{bijection}~; cette
application, étant fonctorielle en $\mathcal{F}$, définira alors un
isomorphisme de foncteurs en $\mathcal{F}$. Le couple
$(\mathcal{G}',\rho)$ étant solution d'un problème universel, on sait qu'il
est \emph{déterminé à un isomorphisme unique près} lorsqu'il existe. On écrira
alors $\mathcal{G}'=\psi^*(\mathcal{G})$, $\rho=\rho_{\mathcal{G}}$, et par
abus de langage, on dira que $\psi^*(\mathcal{G})$ est \emph{le faisceau image
réciproque} de $\mathcal{G}$ par $\psi$, étant entendu que
$\psi^*(\mathcal{G})$ est considéré comme muni du $\psi$-morphisme
\emph{canonique}
$\rho_{\mathcal{G}}:\mathcal{G}\to\psi^*(\mathcal{G})$, c'est-à-dire de
l'\emph{homomorphisme canonique} de préfaisceaux sur $Y$~:
\begin{equation*}
  \rho_{\mathcal{G}}:\mathcal{G}\to\psi_*(\psi^*(\mathcal{G})).
\tag{3.5.3.2}\label{0.3.5.3.2-fr}
\end{equation*}
Pour tout homomorphisme
$v:\psi^*(\mathcal{G})\to\mathcal{F}$ (où $\mathcal{F}$ est un faisceau sur
$X$ à valeurs dans $\mathbf{K}$), on posera
$v^b=\psi_*(v)\circ\rho_{\mathcal{G}}:
\mathcal{G}\to\psi_*(\mathcal{F})$. Par définition, \emph{tout} morphisme de
préfaisceaux $u:\mathcal{G}\to\psi_*(\mathcal{F})$ est de la forme $v^b$ pour
un $v$ et un seul, que l'on notera $u^\sharp$. En d'autres termes, tout
morphisme $u:\mathcal{G}\to\psi_*(\mathcal{F})$ de préfaisceaux se factorise
de façon unique en
\begin{equation*}
  u:\mathcal{G}\xrightarrow{\rho_{\mathcal{G}}}
  \psi_*(\psi^*(\mathcal{G}))
  \xrightarrow{\psi_*(u^\sharp)}\psi_*(\mathcal{F}).
\tag{3.5.3.3}\label{0.3.5.3.3-fr}
\end{equation*}
\end{env}

